% E1.8 -- a rota do intervalo (base de Cohen-Daubechies-Vial), o caso eps > 0
% e a analise de s' dos cenarios de wafc/R/dgp.R.
%
% Este arquivo NAO vai ao manuscrito (D26): o artigo enuncia a teoria na base
% periodizada com eps = 0. Ele existe por dois motivos: deixar pronta, e
% escrita, a rota sem periodicidade para o dia em que um referee a pedir; e
% guardar a analise de regularidade efetiva que E2.4 e E4 vao citar (D27).
%
% Numeracao global dos resultados (docs/TAREFA.md, §5): este arquivo continua
% de Proposicao 4, Lema 10, Teorema 2 e Corolario 5, e imprime o numero global
% via \setcounter, como 04 e 05.
%
% Conferencia numerica: derivations/check/07-rota-intervalo.R (imprime OK).
% Notacao: docs/notacao.md (congelada em E1.1). Macros: derivations/macros.tex.
\documentclass[11pt,a4paper]{article}
\usepackage[T1]{fontenc}
\usepackage[utf8]{inputenc}
\usepackage[brazil]{babel}
\usepackage[margin=2.6cm]{geometry}
\usepackage{enumitem}
\usepackage{booktabs}
\usepackage{xcolor}
\usepackage[colorlinks=true,linkcolor=blue!60!black,citecolor=blue!60!black]{hyperref}
% Macros do WAFC. Implementa docs/notacao.md (esboço; congela em E1.1).
% Único lugar onde uma macro é definida; os derivations/*.tex fazem % Macros do WAFC. Implementa docs/notacao.md (esboço; congela em E1.1).
% Único lugar onde uma macro é definida; os derivations/*.tex fazem % Macros do WAFC. Implementa docs/notacao.md (esboço; congela em E1.1).
% Único lugar onde uma macro é definida; os derivations/*.tex fazem \input{macros}.
\usepackage{amsmath,amssymb,amsthm}
\newcommand{\R}{\mathbb{R}}
\newcommand{\E}{\mathbb{E}}
\newcommand{\Prob}{\mathbb{P}}
\newcommand{\T}{^{\top}}
\newcommand{\norm}[1]{\left\lVert #1 \right\rVert}
\newcommand{\normn}[1]{\left\lVert #1 \right\rVert_{n}}
\newcommand{\normL}[1]{\left\lVert #1 \right\rVert_{L_2(P)}}
\newcommand{\Besov}[3]{B^{#1}_{#2,#3}}
\newcommand{\bX}{\mathbf{X}}
\newcommand{\bU}{\mathbf{U}}
\newcommand{\bZ}{\mathbf{Z}}
\newcommand{\btheta}{\boldsymbol{\theta}}
\newcommand{\thetastar}{\btheta^{*}}
\newcommand{\bc}{\mathbf{c}}
\newcommand{\PiJ}{\Pi_{J}}
\newcommand{\VJ}{V_{J}}
\newcommand{\NJ}{N_{J}}
\newtheorem{proposition}{Proposition}
\newtheorem{lemma}{Lemma}
\newtheorem{theorem}{Theorem}
\newtheorem{corollary}{Corollary}
\theoremstyle{definition}
\newtheorem{assumption}{Assumption}
\newtheorem{remark}{Remark}
.
\usepackage{amsmath,amssymb,amsthm}
\newcommand{\R}{\mathbb{R}}
\newcommand{\E}{\mathbb{E}}
\newcommand{\Prob}{\mathbb{P}}
\newcommand{\T}{^{\top}}
\newcommand{\norm}[1]{\left\lVert #1 \right\rVert}
\newcommand{\normn}[1]{\left\lVert #1 \right\rVert_{n}}
\newcommand{\normL}[1]{\left\lVert #1 \right\rVert_{L_2(P)}}
\newcommand{\Besov}[3]{B^{#1}_{#2,#3}}
\newcommand{\bX}{\mathbf{X}}
\newcommand{\bU}{\mathbf{U}}
\newcommand{\bZ}{\mathbf{Z}}
\newcommand{\btheta}{\boldsymbol{\theta}}
\newcommand{\thetastar}{\btheta^{*}}
\newcommand{\bc}{\mathbf{c}}
\newcommand{\PiJ}{\Pi_{J}}
\newcommand{\VJ}{V_{J}}
\newcommand{\NJ}{N_{J}}
\newtheorem{proposition}{Proposition}
\newtheorem{lemma}{Lemma}
\newtheorem{theorem}{Theorem}
\newtheorem{corollary}{Corollary}
\theoremstyle{definition}
\newtheorem{assumption}{Assumption}
\newtheorem{remark}{Remark}
.
\usepackage{amsmath,amssymb,amsthm}
\newcommand{\R}{\mathbb{R}}
\newcommand{\E}{\mathbb{E}}
\newcommand{\Prob}{\mathbb{P}}
\newcommand{\T}{^{\top}}
\newcommand{\norm}[1]{\left\lVert #1 \right\rVert}
\newcommand{\normn}[1]{\left\lVert #1 \right\rVert_{n}}
\newcommand{\normL}[1]{\left\lVert #1 \right\rVert_{L_2(P)}}
\newcommand{\Besov}[3]{B^{#1}_{#2,#3}}
\newcommand{\bX}{\mathbf{X}}
\newcommand{\bU}{\mathbf{U}}
\newcommand{\bZ}{\mathbf{Z}}
\newcommand{\btheta}{\boldsymbol{\theta}}
\newcommand{\thetastar}{\btheta^{*}}
\newcommand{\bc}{\mathbf{c}}
\newcommand{\PiJ}{\Pi_{J}}
\newcommand{\VJ}{V_{J}}
\newcommand{\NJ}{N_{J}}
\newtheorem{proposition}{Proposition}
\newtheorem{lemma}{Lemma}
\newtheorem{theorem}{Theorem}
\newtheorem{corollary}{Corollary}
\theoremstyle{definition}
\newtheorem{assumption}{Assumption}
\newtheorem{remark}{Remark}

\newcommand{\bPsi}{\boldsymbol{\Psi}}
\newcommand{\bPhi}{\boldsymbol{\Phi}}
\newcommand{\lmin}{\lambda_{\min}}
\newcommand{\lmax}{\lambda_{\max}}
\newcommand{\Amat}{A}
\newcommand{\Bmat}{B}
\newcommand{\PA}{\Pi_{\Amat}}
\newcommand{\bveps}{\boldsymbol{\varepsilon}}
\newcommand{\eps}{\varepsilon}
\newcommand{\Ieps}{I_{\eps}}
\newcommand{\Cpsiint}{C_{\psi}^{\mathrm{int}}}
\newcommand{\dint}{d^{\mathrm{int}}}
\newcommand{\pzero}{p_{0}}
\newcommand{\gtil}{\widetilde{\gamma}}   % como em 04-oraculo.tex

% A numeracao global comeca depois de Proposicao 4, Lema 10, Teorema 2,
% Corolario 5 (docs/TAREFA.md, §5).
\setcounter{proposition}{4}
\setcounter{lemma}{10}
\setcounter{theorem}{2}
\setcounter{corollary}{5}

\title{E1.8. A rota do intervalo, o preço exato da margem e a regularidade
efetiva dos cenários}
\author{wafc-draft, \texttt{derivations/07-rota-intervalo.tex}}
\date{2026-09-19}

\begin{document}
\maketitle

\section*{Onde este arquivo entra}

A decisão D26 fixou que a teoria do artigo fica na base periodizada com
$\eps = 0$, e que a margem do reescalonamento é dispositivo de amostra
finita. A mesma decisão registrou duas rotas de princípio para o caso de um
referee pedir teoria sem periodicidade, e este arquivo escreve a primeira
delas por inteiro, mais o que sobrou por escrever das outras duas frentes.
São três partes independentes:

\begin{description}[nosep]
  \item[Parte A.] A rota do intervalo. O que transfere da base periodizada
    para a base de Cohen, Daubechies e Vial (CDV), com o que é verbatim e o
    que não é, e os dois pontos que não transferem de graça: a cota pontual
    $\sum_{\lev\tsl}\wav{\lev}{\tsl}^{2}(u) \le C_{\psi}2^{J}$, que E1.4
    conferiu só na base periódica, e o efeito dos $pq(2^{\jzero}-1)$
    coeficientes de escala não penalizados no termo
    $\sigma^{2}\times(\text{não penalizados})/n$ do Teorema~1.
  \item[Parte B.] O caso $\eps > 0$, consolidando o Lema~10 e D26: o
    cruzamento exato das duas exigências que a margem tem de satisfazer, por
    que o ganho assintótico é exatamente zero, e qual enunciado é verdadeiro
    (um resultado de amostra finita, e não um teorema de taxa).
  \item[Parte C.] A regularidade efetiva $\seff$ de cada componente dos
    cenários de \texttt{wafc/R/dgp.R}, nos três regimes, que é o que D27
    declara e o que E2.4 e E4 vão citar.
\end{description}

Nada aqui altera o manuscrito, e nada aqui altera arquivo fechado de E1: as
Partes A e B citam E1.2 a E1.6 e E1.3b sem mexer neles.

\section*{Convenções}

A notação é a de \texttt{docs/notacao.md}. $L$ é o comprimento do filtro
(\texttt{filter.size}), $N = L/2$ o número de momentos nulos para Daublets e
Symmlets, e $\operatorname{supp}\psi \subseteq [0, L-1]$. As constantes
$B_{X}$, $c_{U}$, $C_{U}$, $\kappa_{1}$, $\kappa_{2}$ são as das Hipóteses de
E1.3 e E1.4. Escrevemos $\seff = s - (1/\pi - 1/2)_{+}$, e ``queda por
nível'' para $-\log_{2}\bigl(\norm{\theta_{\lev+1,\cdot}}_{2} /
\norm{\theta_{\lev\cdot}}_{2}\bigr)$, que é $\seff$ quando
$\norm{\theta_{\lev\cdot}}_{2} \asymp 2^{-\lev\seff}$: é a grandeza que as
três conferências numéricas do projeto medem.

Os três regimes de borda que a Parte~C compara, e que a Parte~B separa, são:

\begin{description}[nosep]
  \item[(a) periódico com $\eps = 0$.] A base é a periodizada de D2, a
    moduladora vive em todo $[0,1]$, e a hipótese de regularidade é a
    Hipótese~2 de E1.3, em forma de sequência. É o regime do manuscrito.
  \item[(b) periódico com margem e extensão.] A base é a mesma, mas
    $\operatorname{supp}P_{U_{\md}} \subseteq \Ieps = [\eps, 1-\eps]$ e a
    componente é trocada pela extensão do Lema~10 (D23). É o regime em que o
    código roda, com \texttt{rescale = TRUE} e $\eps$ fixo.
  \item[(c) intervalo.] A base é a CDV, não há periodicidade, não há margem e
    não há extensão. É a rota da Parte~A.
\end{description}

\part*{Parte A. A rota do intervalo}
\addcontentsline{toc}{part}{Parte A}

\section{A base CDV e a reparametrização}

A base de Cohen, Daubechies e Vial (1993) mantém, em cada nível $\lev$, os
$2^{\lev} - L$ transladados interiores $\wav{\lev}{\tsl}(u) =
2^{\lev/2}\psi(2^{\lev}u - \tsl)$ e os completa com $L/2$ funções de borda em
cada extremo, corrigidas de modo que $V_{J}$ reproduza os polinômios de grau
$< N$ em $[0,1]$ e que o sistema
$\{\phi^{\mathrm{int}}_{\jzero\tsl}\}_{\tsl<2^{\jzero}} \cup
\{\psi^{\mathrm{int}}_{\lev\tsl}\}_{\jzero\le\lev<J}$ seja ortonormal em
$L_{2}[0,1]$ com a medida de Lebesgue. A construção exige
$2^{\jzero}\gtrsim 2L$ para que as duas zonas de borda não se toquem; o
\texttt{wbasis()} impõe $\jzero \gtrsim \log_{2}(5L)$, isto é
$\jzero \ge 4$ com $L = 8$ e $\jzero \ge 6$ com o $L = 20$ do \texttt{wall()}.

Duas diferenças em relação ao caso periódico organizam tudo o que segue.
Primeira: a constante pertence a $V_{\jzero}$, mas não é uma coluna. Como
$\mathbf{1}$ é polinômio de grau zero, $\mathbf{1} = \sum_{\tsl}\mu_{\tsl}
\phi^{\mathrm{int}}_{\jzero\tsl}$ com $\mu_{\tsl} =
\int_{0}^{1}\phi^{\mathrm{int}}_{\jzero\tsl}$ e $\norm{\mu}_{2} = 1$; não há
uma única coluna a descartar, como em D2. Segunda: as wavelets continuam
centradas, porque $\psi^{\mathrm{int}}_{\lev\tsl} \perp V_{\jzero} \ni
\mathbf{1}$, mas as funções de escala não.

A §5 de \texttt{01-identificabilidade.md} já resolveu isso, e a solução é a
que este arquivo usa em toda a Parte~A. Seja $Q \in
\R^{2^{\jzero}\times(2^{\jzero}-1)}$ uma base ortonormal de
$\{\alpha : \mu\T\alpha = 0\}$ e $\bPhi(u) =
(\phi^{\mathrm{int}}_{\jzero\tsl}(u))_{\tsl}$. O bloco $(\cv,\md)$ passa a
ter as colunas
\begin{equation}
  \label{eq:bloco}
  \bigl[\,X_{\cv}\,\bPhi(U_{\md})\T Q \;\bigm|\;
  X_{\cv}\,\bPsi^{\mathrm{int}}(U_{\md})\T\,\bigr],
\end{equation}
com $2^{\jzero}-1$ colunas de escala, \emph{não penalizadas}, e
$2^{J} - 2^{\jzero}$ wavelets, penalizadas; o total, $2^{J}-1$, é o mesmo
$\NJ$ do caso periódico. O que importa para tudo o que segue é que as colunas
de $\bPhi\T Q$ são ortonormais em $L_{2}[0,1]$ ($Q\T Q = I$) e têm integral
zero ($\mu\T Q = 0$). Escrevemos
\[
  \dint = pq\,(2^{J} - 2^{\jzero}), \qquad
  \pzero = p\bigl\{1 + q\,(2^{\jzero}-1)\bigr\}
\]
para a dimensão penalizada e para o número de colunas não penalizadas,
e $\Amat^{\mathrm{int}}$, $\Bmat^{\mathrm{int}}$ para os dois blocos da
matriz de desenho, na partição do Lema~4 de E1.5.

\section{O que transfere verbatim}

\paragraph{Aproximação (E1.3).} O Lema~2, o Lema~3 e o Corolário~1 valem com
$\Besov{s}{\pi}{r}[0,1]$, o espaço de Besov do intervalo \emph{sem condição de
periodicidade}, no lugar da classe periódica, e com a mesma prova. A razão é
que a base CDV caracteriza $\Besov{s}{\pi}{r}[0,1]$ pelos coeficientes (Cohen,
2003, \S 3.9), de modo que a Hipótese~2 de E1.3, em forma de sequência, passa
a ser consequência de pertencimento a $\Besov{s}{\pi}{r}[0,1]$, e não uma
hipótese periódica disfarçada. Os quatro passos da prova do Lema~2 são
insensíveis à troca: o Passo~1 usa completude e Parseval, o Passo~2 é a
desigualdade de Hölder dentro do nível, o Passo~3 é a soma geométrica e o
Passo~4 é a densidade limitada. O Lema~3 usa a sobreposição limitada, que vale
com a constante ajustada (Lema~\ref{lem:pontual} abaixo). E a projeção de uma
função de média zero em $V_{J}$ continua ortogonal a $\mathbf{1}$, porque
$\mathbf{1}\in V_{\jzero}\subseteq V_{J}$.

O que desaparece é o que E1.3b teve de construir: não há operador de
extensão, não há corte $\chi_{\eps}$, não há a constante
$\tilde C_{g} \asymp \eps^{-(s-1/\pi)}$ e não há o termo $2^{-J}$ da
Proposição~2. A Proposição~2 simplesmente não tem análogo aqui, porque a
emenda $0 \equiv 1$ não existe.

\paragraph{Identificabilidade (E1.2).} Já está feita, na §5 de
\texttt{01-identificabilidade.md}. A Proposição~1 não toca na base. O Lema~1
vale literalmente com $[\bPhi Q \mid \bPsi^{\mathrm{int}}]$ no lugar de
$\bPsi$, porque a prova do Lema~1(ii) usa exatamente duas propriedades das
colunas do bloco, ortonormalidade em $L_{2}[0,1]$ e integral zero, e a
reparametrização \eqref{eq:bloco} entrega as duas. Manter as $2^{\jzero}$
colunas de escala sem a restrição repete a constante em cada bloco e deixa
$\Gram$ singular com nulidade $pq$, que é o que o \texttt{wall()} faz hoje com
\texttt{boundary = "interval"}.

\paragraph{Desenho (E1.4).} A Proposição~3 vale verbatim, com $\Cpsiint$ no
lugar de $C_{\psi}$. A prova da parte~(i), que é onde a base entra, usa
apenas que $\{\mathbf{1}\}\cup\{\text{colunas do bloco}\}$ é ortonormal em
$L_{2}([0,1]^{q}, du)$: dentro de uma coordenada isso é a ortonormalidade da
base CDV mais $Q\T Q = I$; entre coordenadas, é o fato de toda coluna do
bloco ter integral zero, que é $\mu\T Q = 0$ para as de escala e
$\psi^{\mathrm{int}}\perp\mathbf{1}$ para as wavelets. Logo
$\int_{[0,1]^{q}}h_{a}(u)^{2}du = \norm{a}_{2}^{2}$ e
\[
  c_{U}\le\lmin(\Gram_{\bPsi})\le\lmax(\Gram_{\bPsi})\le C_{U},
  \qquad
  \kappa_{1}c_{U} \;\le\; \lmin(\Gram) \;\le\; \lmax(\Gram) \;\le\;
  \kappa_{2}C_{U},
\]
sem $\lmin(G_{\eps})$ e sem hipótese de suporte próprio. As partes (ii) e
(iii) não mencionam a base; a parte (iv) precisa de
$R_{J}^{\mathrm{int}} = pB_{X}^{2}(1 + q\Cpsiint 2^{J})$, e é o
Lema~\ref{lem:pontual} que a sustenta.

\paragraph{Oráculo e taxas (E1.5, E1.6).} Consomem só essas peças. O Lema~4
(perfilagem) é álgebra de projeção e não conhece a base; exige apenas que
$\Amat^{\mathrm{int}}$ tenha posto $\pzero$. O Lema~5 (calibração) é a
desigualdade máxima sub-gaussiana sobre as $\dint$ colunas penalizadas. O
Lema~6 usa $|X_{\cv}|\le B_{X}$ e a densidade marginal limitada, e vale com
$\norm{\psi}_{\infty}$ trocado pelo máximo entre $\norm{\psi}_{\infty}$ e as
normas das $L$ funções de borda. O Lema~7 é a Hipótese de Besov e a soma
geométrica, agora de $\lev = \jzero$ a $J-1$, o que só pode diminuir
$\norm{\thetastar}_{1}$. O Teorema~1 sai desses quatro. E o Teorema~2, a
Proposição~4 e os Corolários~4 e~5 de E1.6 consomem o Teorema~1 e a
Proposição~3.

\section{O que não transfere de graça}

\subsection*{A cota pontual}

A Hipótese~3 de E1.4 pede uma constante $C_{\psi}$ com
$\sup_{u}\sum_{\lev<J}\sum_{\tsl}\wav{\lev}{\tsl}(u)^{2} \le C_{\psi}2^{J}$, e
a justificativa dada lá (no máximo $L$ wavelets periodizadas não nulas em cada
$u$, cada uma de norma $2^{\lev/2}\norm{\psi}_{\infty}$) é específica da base
periódica: na CDV as funções de borda não são transladadas de $\psi$.

\begin{lemma}[cota pontual na base do intervalo]\label{lem:pontual}
Seja a base CDV de $V_{J}$ em $[0,1]$, com $\jzero$ admissível e $L$ o
comprimento do filtro, e seja $W(u) \in \R^{2^{J}}$ o vetor de todas as suas
funções avaliadas em $u$. Então existe $\Cpsiint < \infty$, que depende apenas
da família de wavelets e de $\jzero$, tal que
\[
  \sup_{u\in[0,1]}\norm{W(u)}_{2}^{2}
  \;=\; \sup_{u}\Bigl\{\sum_{\tsl}\phi^{\mathrm{int}}_{\jzero\tsl}(u)^{2}
  + \sum_{\lev=\jzero}^{J-1}\sum_{\tsl}\psi^{\mathrm{int}}_{\lev\tsl}(u)^{2}
  \Bigr\}
  \;\le\; \Cpsiint\,2^{J}
  \qquad\text{para todo } J > \jzero .
\]
A mesma cota, com a mesma constante, vale para o bloco reparametrizado
$[\bPhi Q \mid \bPsi^{\mathrm{int}}]$ de \eqref{eq:bloco}.
\end{lemma}

\begin{proof}
Fixe $\lev \ge \jzero$. As funções do nível $\lev$ são de três tipos. As
$2^{\lev}-L$ \emph{interiores} são $2^{\lev/2}\psi(2^{\lev}u - \tsl)$; como
$\operatorname{supp}\psi\subseteq[0,L-1]$, no máximo $L-1$ delas são não nulas
em cada $u$, e cada uma é $\le 2^{\lev/2}\norm{\psi}_{\infty}$. As $L/2$
\emph{de borda à esquerda} são $2^{\lev/2}\psi^{\mathrm{esq}}_{r}(2^{\lev}u)$,
$r < L/2$, com $\psi^{\mathrm{esq}}_{r}$ um conjunto \emph{fixo} de funções
limitadas de suporte compacto em $[0, C_{L}]$, independentes de $\lev$ e de
$J$ (é a construção de Cohen, Daubechies e Vial, que corrige as wavelets da
borda uma vez e as dilata depois); as de borda à direita são a imagem delas
por $u\mapsto 1-u$. Em cada $u$, no máximo $L/2$ de cada lado são não nulas.
Logo
\[
  \sum_{\tsl}\psi^{\mathrm{int}}_{\lev\tsl}(u)^{2}
  \;\le\; \Bigl\{(L-1)\norm{\psi}_{\infty}^{2}
  + L\max_{r}\norm{\psi^{\mathrm{esq}}_{r}}_{\infty}^{2}\Bigr\}\,2^{\lev}
  \;=:\; c_{1}\,2^{\lev},
\]
e $\sum_{\lev=\jzero}^{J-1}c_{1}2^{\lev} \le c_{1}2^{J}$. O bloco de escala
obedece ao mesmo argumento com $\phi$ e as suas funções de borda, dando
$\sum_{\tsl}\phi^{\mathrm{int}}_{\jzero\tsl}(u)^{2} \le c_{2}2^{\jzero} \le
c_{2}2^{J}$. Tome $\Cpsiint = c_{1} + c_{2}$.

Para o bloco reparametrizado, escreva $\widetilde W(u) = (Q\T\bPhi(u),
\bPsi^{\mathrm{int}}(u))$. Como $QQ\T$ é a projeção ortogonal sobre
$\{\mu\}^{\perp}$, $\norm{Q\T\bPhi(u)}_{2}^{2} \le \norm{\bPhi(u)}_{2}^{2}$, e
portanto $\norm{\widetilde W(u)}_{2}^{2}\le\norm{W(u)}_{2}^{2}$.
\end{proof}

A conferência numérica mede $\Cpsiint$ e encontra $7{,}55$ (estável em
$J = 4,\dots,9$, com leve decrescimento até $6{,}85$), contra $1{,}30$ a
$1{,}36$ na base periódica: um fator $5{,}6$. A constante entra linearmente em
$R_{J}$, logo entra linearmente na condição empírica de E1.4(iv): a rota do
intervalo exige, para a mesma probabilidade, cerca de seis vezes o $n$ que a
periódica exige. É um preço de constante, não de taxa.

\subsection*{Os coeficientes de escala não penalizados}

\begin{proposition}[a rota do intervalo, enunciada]\label{prop:cdv}
Considere o desenho \eqref{eq:bloco} com a base CDV de nível inicial
$\jzero$, valendo as Hipóteses~1 e~2 de E1.5, as Hipóteses de E1.4 com
$\Cpsiint$ do Lema~\ref{lem:pontual} no lugar de $C_{\psi}$, e a Hipótese de
Besov de E1.3 lida em $\Besov{s}{\pi}{r}[0,1]$, \emph{sem periodicidade,
sem margem e sem extensão}. Suponha $\Amat^{\mathrm{int}}$ de posto $\pzero$.
Então:
\begin{enumerate}[label=\textup{(\roman*)}, leftmargin=*]
  \item a Proposição~3 de E1.4 vale como está, com
    $\lmin(\Gram) \ge \kappa_{1}c_{U}$,
    $\lmax(\Gram) \le \kappa_{2}C_{U}$ e
    $R_{J}^{\mathrm{int}} = pB_{X}^{2}(1+q\Cpsiint 2^{J})$;
  \item o Corolário~1 de E1.3 vale com a mesma cota de viés
    $\mathcal{B}_{J}^{2}\asymp (pq)^{2}B_{X}^{2}C_{U}C_{g}^{2}2^{-2J\seff}$,
    sem o termo $2^{-J}$ e sem o fator $\eps^{-2(s-1/\pi)}$;
  \item o Teorema~1 de E1.5 vale com $\Amat$ trocado por
    $\Amat^{\mathrm{int}}$, $s_{0}$ contado entre as $\dint$ colunas
    penalizadas, e o termo dos níveis trocado por
    \[
      \E\bigl[\normn{\PA\bveps}^{2}\bigm|\bX,\bU\bigr]
      \;\le\; \frac{\sigma^{2}\,\pzero}{n}
      \;=\; \frac{\sigma^{2}\,p\{1+q(2^{\jzero}-1)\}}{n};
    \]
  \item o Teorema~2 e os Corolários~4 e~5 de E1.6 valem com as mesmas taxas,
    porque $\pzero$ não depende de $n$ e $\sigma^{2}\pzero/n$ é de ordem
    estritamente menor que $(J_{n}/n)^{2\seff/(2\seff+1)}$.
\end{enumerate}
\end{proposition}

\begin{proof}
(i) é a §2 acima mais o Lema~\ref{lem:pontual}: a parte~(i) da Proposição~3 só
usa ortonormalidade em $L_{2}([0,1]^{q},du)$, que \eqref{eq:bloco} entrega, e
a parte~(iv) só usa $\norm{\bZ_{i}}_{2}^{2}\le R_{J}$, que o
Lema~\ref{lem:pontual} dá. (ii) é a §2, parágrafo de aproximação. (iii): a
prova do Teorema~1 nunca usa o número de colunas de $\Amat$ além de dois
lugares, o posto (Lema~4) e a cota
$\E[\normn{\PA\bveps}^{2}\mid\bX,\bU] = \sigma^{2}\operatorname{rank}
(\Amat)/n$ da parte~(iii); trocar $p$ por $\pzero$ é trocar o posto. Note que
as colunas de escala entram em $\Amat^{\mathrm{int}}$, e não em
$\Bmat^{\mathrm{int}}$, porque não são penalizadas; a perfilagem do Lema~4 as
remove do problema exatamente como remove os $\lvl{\cv}$, e
$\gtil \ge \lmin(\Gramhat)$ continua valendo pelo complemento de Schur.
(iv): no Teorema~2, $2\seff/(2\seff+1) < 1$, e $\pzero$ é constante em $n$.
\end{proof}

\begin{corollary}[as taxas sem periodicidade]\label{cor:taxa-cdv}
Nas condições da Proposição~\ref{prop:cdv}, com
$\penn$ e $J_{n}$ do Teorema~2 de E1.6,
\[
  \normn{\widehat f - f}^{2}
  \;=\; O_{p}\!\left(\Bigl(\frac{\log n}{n}\Bigr)^{2\seff/(2\seff+1)}\right),
  \qquad\text{e, pelo Corolário~5,}\qquad
  O_{p}\!\left(\Bigl(\frac{\log n}{n}\Bigr)^{2s/(2s+1)}\right),
\]
com $\gcomp{\cv}{\md}\in\Besov{s}{\pi}{r}[0,1]$ e nenhuma hipótese de
periodicidade, de margem ou de suporte próprio.
\end{corollary}

\subsection*{O que a rota custa}

O preço é inteiramente de amostra finita, e é de dois tipos.

\begin{center}
\begin{tabular}{lll}
  \toprule
  & periódica (D2) & intervalo (CDV) \\
  \midrule
  nível inicial $\jzero$ & $0$ & $\ge 4$ ($L=8$), $\ge 6$ ($L=20$) \\
  colunas não penalizadas $\pzero$ & $p$ & $p\{1+q(2^{\jzero}-1)\}$ \\
  \quad com $p=3$, $q=2$, $L=8$ & $3$ & $93$ \\
  constante pontual $C_{\psi}$ & $1{,}30$ a $1{,}36$ & $\approx 7{,}55$ \\
  $J$ acessíveis & todos & $J > \jzero$ \\
  hipótese de regularidade & periódica & $\Besov{s}{\pi}{r}[0,1]$ \\
  \bottomrule
\end{tabular}
\end{center}

O termo $\sigma^{2}\pzero/n$ é $0{,}093\,\sigma^{2}$ em $n = 1000$ e
$0{,}372\,\sigma^{2}$ em $n = 250$, contra $0{,}003\,\sigma^{2}$ e
$0{,}012\,\sigma^{2}$ no desenho periódico. A conferência numérica mede
$\E\normn{\PA\bveps}^{2} = 0{,}0922$ contra os $0{,}0930$ previstos em
$n = 1000$, $\sigma = 1$: a cota é exata em média, como a Observação do
Teorema~1 já dizia. Nos $n$ de E2.4 a rota do intervalo gasta, só em
coeficientes de escala, uma fração não desprezível do orçamento de erro; é
por isso que ela é a saída de recurso e não o padrão, ainda que
assintoticamente não custe nada. O segundo tipo de preço, $\jzero \ge 4$,
tira $J \in \{1,2,3,4\}$ da grade de \texttt{cv.wafc()}, isto é, exatamente
a faixa que E2.3 encontrou como ótima nos $n$ do piloto.

\part*{Parte B. O caso $\eps > 0$}
\addcontentsline{toc}{part}{Parte B}

\section{As duas exigências, e por que não se encontram}

Sob a Hipótese~5 de E1.3b a moduladora vive em $\Ieps = [\eps, 1-\eps]$ e a
base é periodizada no intervalo maior. Duas coisas dependem de $\eps$ e de
$J$ em direções opostas.

\emph{A Gram restrita.} A Observação do Lema~10 sobre E1.4 mostra que a
constante de desenho passa a ser $c_{U}\lmin(G_{\eps})$, com
$G_{\eps} = \int_{\Ieps^{q}}\bPsi(u)\bPsi(u)\T du$, e que
$\lmin(G_{\eps}) = 0$ assim que $\VJ$ contém função não nula suportada na
margem.

\emph{A faixa contaminada.} O que a margem compra é que as wavelets que
atravessam a emenda $0\equiv 1$, que são as responsáveis pelo termo $2^{-J}$
da Proposição~2, deixem de ver a função. Isso acontece quando o suporte delas
cabe dentro da margem.

As duas coisas são a mesma conta, na mesma escala, e é isso que as torna
incompatíveis.

\begin{proposition}[o cruzamento exato]\label{prop:cruzamento}
Base periodizada de uma família de suporte $[0,L-1]$,
$\operatorname{supp}P_{U_{\md}}\subseteq\Ieps$ com $\eps\in(0,1/4]$, e
$T(\eps) = (L-1)/(2\eps)$. Na circunferência, a margem é um único arco de
comprimento $2\eps$ centrado na emenda. Então:
\begin{enumerate}[label=\textup{(\roman*)}, leftmargin=*]
  \item \textup{(Gram)} se $2^{J} \ge T(\eps)$, então $\lmin(G_{\eps}) = 0$;
    logo $\lmin(G_{\eps}) > 0$ \emph{exige} $2^{J} < T(\eps)$;
  \item \textup{(faixa contaminada)} se $2^{J} \ge 2\,T(\eps)$, então toda
    wavelet de todo nível $\lev\ge J$ cujo suporte contém a emenda está
    contida na margem, e portanto tem coeficiente nulo contra qualquer função
    suportada em $\Ieps$: a cauda $\lev\ge J$ fica livre do termo de
    periodização \emph{sem invocar o Lema~10}, e portanto sem a constante
    $\tilde C_{g}$;
  \item \textup{(disjunção)} as faixas de (i) e (ii) são disjuntas e
    separadas por exatamente um nível de resolução, pois
    $2T(\eps) > T(\eps)$: nenhum par $(J,\eps)$ satisfaz as duas;
  \item \textup{(assintótico)} manter $\lmin(G_{\eps}) > 0$ ao longo de uma
    sequência $J_{n}\to\infty$ obriga $\eps_{n} < (L-1)2^{-J_{n}}/2$, isto é
    $\eps_{n}\asymp 2^{-J_{n}}$; pela Observação do Lema~10 sobre margens que
    encolhem, o expoente garantido nessa escala é $\seff - (s-1/\pi)$, que
    para $\pi = 2$ é exatamente $1/2$, o mesmo que a Proposição~2 entrega
    \emph{sem margem nenhuma}. O ganho assintótico da margem é exatamente
    zero.
\end{enumerate}
\end{proposition}

\begin{proof}
(i) $\VJ$ é gerado pelas funções de escala $\phi_{J\tsl}$, de suporte de
comprimento $(L-1)2^{-J}$. Se $(L-1)2^{-J}\le 2\eps$, alguma delas cabe
inteira no arco da margem, é não nula e está no núcleo de $G_{\eps}$, pois
$G_{\eps}$ integra só sobre $\Ieps$. A condição é $2^{J}\ge(L-1)/(2\eps) =
T(\eps)$.

(ii) Uma wavelet de nível $\lev$ cujo suporte contém a emenda está suportada
num arco de comprimento $(L-1)2^{-\lev}$ que contém o ponto de emenda; a
margem se estende $\eps$ para cada lado dele. Qualquer que seja a fase, o
arco está contido na margem desde que $(L-1)2^{-\lev}\le\eps$, isto é
$2^{\lev}\ge(L-1)/\eps = 2T(\eps)$. Como $\lev\ge J$ e $2^{J}\ge 2T(\eps)$,
vale em todo nível da cauda. Nessa situação o coeficiente contra uma função
que se anula na margem é zero, e a separação entre interiores e borda feita na
prova da Proposição~2 não tem segunda parcela.

(iii) é $2T(\eps) > T(\eps)$.

(iv) Por (i), $2^{J_{n}} < T(\eps_{n})$, isto é
$\eps_{n} < (L-1)2^{-J_{n}}/2$. Escrevendo $\eps_{n} = 2^{-cJ_{n}}$, isso
força $c \ge 1$ a menos de constante; a Observação do Lema~10 dá, para
$\eps = 2^{-cJ}$, expoente garantido $\seff - c(s-1/\pi)$, que em $c = 1$ e
$\pi = 2$ vale $s - (s-1/2) = 1/2$.
\end{proof}

\begin{remark}[por que o zero não é acidente]\label{rem:zero}
As duas exigências medem a mesma largura, $(L-1)2^{-J}$, que é o diâmetro do
suporte de uma função da base no nível mais fino. Excluir a região
contaminada pede que a margem seja \emph{mais larga} que esse diâmetro;
preservar a Gram pede que seja \emph{mais estreita} que ele. A margem teria de
ser as duas coisas. O $1{,}9^{-J}$ herdado do \texttt{wall()} fica logo acima
do cruzamento, e é por isso que E1.3b mediu ali um ganho \emph{parcial}, de
$0{,}96$ bit por nível, e não zero nem a taxa cheia.

A conferência numérica exibe o cruzamento na família
$\eps = a\,(L-1)2^{-J}$, com $J = 5,\dots,9$ e $L - 1 = 7$. Em $a = 0{,}25$ a
Gram vive, com $\lmin(G_{\eps}) = 3{,}5\times 10^{-4}$ estável em $J$, e a
queda por nível do erro restrito é $0{,}50$, isto é, a mesma da periodização
sem margem. Em $a = 0{,}5$ e $a = 1$, $\lmin(G_{\eps})$ é
$1{,}3\times10^{-14}$ e nulo, e o erro restrito desaba para
$1{,}3\times10^{-7}$ e $7{,}5\times10^{-8}$, contra
$2{,}3\times10^{-2}$ em $\eps = 0$ e $2{,}3\times10^{-4}$ em $a = 0{,}25$
(em $J = 9$). O ganho e a degenerescência são o mesmo fenômeno medido de
dois lados: o que some do erro são as direções que somem da Gram.
\end{remark}

\section{Qual enunciado é verdadeiro}

O que sobra depois da Proposição~\ref{prop:cruzamento} não é um teorema de
taxa, é um resultado de amostra finita: fixa-se $\eps$, verifica-se que $J$
está na faixa admissível, e tudo o que E1.5 entrega vale ali, com as
constantes que $\eps$ produz.

\begin{corollary}[o enunciado de amostra finita que a margem
sustenta]\label{cor:amostra-finita}
Fixe $\eps\in(0,1/4]$ e $J$ com $2^{J} < T(\eps)$ e $\lmin(G_{\eps}) > 0$.
Sob as Hipóteses~1 e~5 de E1.3b, as de E1.4 e as de E1.5, o Teorema~1 e os
Corolários~2 e~3 valem com
\[
  \gamma \;=\; \kappa_{1}\,c_{U}\,\lmin(G_{\eps})
  \qquad\text{e}\qquad
  \mathcal{B}_{J}^{2}\;\asymp\;(pq)^{2}B_{X}^{2}C_{U}\,\tilde C_{g}^{2}\,
  2^{-2J\seff}, \quad
  \tilde C_{g}\asymp\eps^{-(s-1/\pi)} ,
\]
$\tilde C_{g}$ do Lema~10(ii). Não há enunciado de taxa associado: para
$\eps$ fixo a faixa admissível de $J$ é finita, e todo $J$ nela é $O(1)$ em
$n$.
\end{corollary}

Três consequências práticas, que são o que E2.4 leva daqui.

\begin{enumerate}[label=\textup{(\alph*)}, leftmargin=*]
  \item \textbf{A faixa admissível é curta e depende de $\eps$.} Com $L = 8$,
    o limiar suficiente de degenerescência é $2^{J}\ge T(\eps) = 7/(2\eps)$,
    isto é $J \ge 8, 7, 6$ para $\eps = 0{,}02,\ 0{,}05,\ 0{,}10$. A
    conferência mede o maior $J$ com $\lmin(G_{\eps})$ acima de $10^{-8}$ e
    encontra $7$, $5$ e $4$: o limiar suficiente não é justo, e a
    degenerescência chega até um nível antes dele.
  \item \textbf{O padrão atual não é binding nos $n$ do piloto.} Com
    $\eps = 0{,}05$ a faixa é $J\le 5$, e E2.3 encontrou o ótimo da validação
    cruzada em $J = 3$ a $5$ para $n\le 1000$. A margem provisória de E2.1b
    está, por sorte ou por escala, dentro da faixa; é isso que E2.4 tem de
    confirmar ao varrer $\eps$, e é o que impede uma escolha de $\eps$ pequeno
    demais de ser gratuita.
  \item \textbf{A escolha de $\eps$ é um compromisso sem ótimo assintótico.}
    Aumentar $\eps$ baixa o viés de borda e encurta a faixa de $J$; diminuir
    $\eps$ alarga a faixa e devolve o viés. A família $a\,(L-1)2^{-J}$ com $a$
    em torno de $1$, que a pergunta~6 do \texttt{ESTADO.md} manda medir, é
    exatamente a que atravessa o cruzamento, e por isso é a família certa para
    localizar o compromisso em ISE, ainda que nenhuma escolha dentro dela
    compre taxa.
\end{enumerate}

\part*{Parte C. A regularidade efetiva dos cenários}
\addcontentsline{toc}{part}{Parte C}

\section{O que fixa $\seff$}

A regularidade efetiva que a aproximação linear enxerga é decidida por duas
coisas, e só por elas: a singularidade mais forte que a função apresenta
dentro da região que a base tem de representar, e o número $N$ de momentos
nulos, que é um teto. O lema abaixo é a forma do primeiro fato que a Parte~C usa; o
segundo é a observação que o segue.

\begin{lemma}[a ordem da descontinuidade fixa a queda por
nível]\label{lem:ordem}
Seja $\psi$ limitada, de suporte $[0,L-1]$, com $N$ momentos nulos. Seja
$h:\R\to\R$ que, numa vizinhança de $t$, é $C^{N}$ de cada lado, com
$h^{(a)}(t^{+}) = h^{(a)}(t^{-})$ para $a = 0,\dots,r-1$ e
$\Delta_{r} := h^{(r)}(t^{+}) - h^{(r)}(t^{-}) \ne 0$, para algum inteiro
$0\le r< N$. Então, para $\lev$ grande:
\begin{enumerate}[label=\textup{(\roman*)}, leftmargin=*, nosep]
  \item os coeficientes $\inner{h}{\psi_{\lev\tsl}}$ das $O(1)$ translações
    cujo suporte contém $t$ são
    \[
      \inner{h}{\psi_{\lev\tsl}} \;=\;
      \frac{\Delta_{r}}{r!}\,2^{-\lev(r+1/2)}\,
      \Lambda_{r}(2^{\lev}t - \tsl) \;+\; O\bigl(2^{-\lev(r+3/2)}\bigr),
      \qquad
      \Lambda_{r}(\tau) = \int_{z>\tau}(z-\tau)^{r}\psi(z)\,dz,
    \]
    e $\Lambda_{r}$ não é identicamente nula em $[0,L-1]$;
  \item os demais coeficientes do nível, dentro da vizinhança, são
    $O(2^{-\lev(N+1/2)})$;
  \item logo $\norm{\theta_{\lev\cdot}}_{2}\asymp 2^{-\lev(r+1/2)}$ e a
    regularidade efetiva imposta pela singularidade é $\seff = r + 1/2$: um
    salto de valor ($r = 0$) dá $1/2$, uma quina ($r = 1$) dá $3/2$.
\end{enumerate}
Se, ao contrário, $h$ é $C^{\infty}$ numa vizinhança de toda a região
representada, então $\norm{\theta_{\lev\cdot}}_{2} = O(2^{-\lev N})$: são
$2^{\lev}$ coeficientes de ordem $2^{-\lev(N+1/2)}$, e a queda por nível é
$N$, o número de momentos nulos, e não a suavidade de $h$.
\end{lemma}

\begin{proof}
Seja $h_{-}$ a extensão $C^{N}$ da peça à esquerda e $H = h - h_{-}$, que se
anula à esquerda de $t$ e vale $(\Delta_{r}/r!)(y-t)^{r} + O((y-t)^{r+1})$ à
direita. Para $h_{-}$, a fórmula de Taylor de ordem $N$ e os $N$ momentos
nulos dão $\inner{h_{-}}{\psi_{\lev\tsl}} = O(2^{-\lev(N+1/2)})$, que é (ii) e
a segunda parcela de (i). Para $H$, a mudança de variável $y = 2^{-\lev}(z +
\tsl)$ dá
\[
  \int_{y>t}(y-t)^{r}\psi_{\lev\tsl}(y)\,dy
  \;=\; 2^{-\lev(r+1/2)}\int_{z>\tau}(z-\tau)^{r}\psi(z)\,dz,
  \qquad \tau = 2^{\lev}t - \tsl ,
\]
e o resto $O((y-t)^{r+1})$ contribui $O(2^{-\lev(r+3/2)})$. Se
$\Lambda_{r}\equiv 0$ em $[0,L-1]$, então $\psi$ seria ortogonal a todas as
potências truncadas $(z-\tau)_{+}^{r}$, o que derivando $r+1$ vezes em $\tau$
dá $\psi \equiv 0$ no suporte. O número de $\tsl$ com $t$ no suporte de
$\psi_{\lev\tsl}$ é no máximo $L-1$, o que é (iii) junto com (ii). No caso
$C^{\infty}$, (ii) vale em todas as $2^{\lev}$ translações e
$\norm{\theta_{\lev\cdot}}_{2}\le 2^{\lev/2}\max_{\tsl}|\theta_{\lev\tsl}| =
O(2^{-\lev N})$.
\end{proof}

Duas leituras do lema, que é o que organiza a tabela da próxima seção. A
primeira: na base periodizada, o ponto de emenda $0\equiv 1$ é um ponto
interior da reta para a extensão $1$-periódica $\tilde g$, e
$\Delta_{r} = g^{(r)}(0) - g^{(r)}(1)$. É por isso que uma função perfeitamente
suave em $[0,1]$ pode ler $\seff = 1/2$: o que a base vê não é $g$, é
$\tilde g$. A segunda: a margem com extensão (regime (b)) e a base do
intervalo (regime (c)) removem esse ponto, a primeira porque troca $g$ por
uma função que se anula perto da emenda, a segunda porque não tem emenda; e
em ambos os casos o que resta é a pior singularidade \emph{interior}, que
nenhuma escolha de borda alcança. O teto $N$ vale nos três regimes.

\section{Os seis componentes de \texttt{dgp.R}}

A conferência numérica mede o valor e a derivada de cada componente nos dois
extremos, que é o que decide $r$ na emenda. O resultado, com
$\Delta$ medido contra a escala da própria função:

\begin{center}
\begin{tabular}{lrrrrl}
  \toprule
  componente & $g(0)$ & $g(1)$ & $g'(0)$ & $g'(1)$ & singularidade que manda \\
  \midrule
  \texttt{sine} $\sin(2\pi u)$      & $0$ & $0$ & $2\pi$ & $2\pi$ & nenhuma: $1$-periódica e $C^{\infty}$ \\
  \texttt{cosine} $\cos(4\pi u)$    & $1$ & $1$ & $0$ & $0$ & nenhuma: $1$-periódica e $C^{\infty}$ \\
  \texttt{cubic} $u^{3}-1{,}4u^{2}+0{,}4u$ & $0$ & $0$ & $0{,}4$ & $0{,}6$ & quina na emenda, $r = 1$ \\
  \texttt{bumps}     & $\approx 0$ & $\approx 0$ & $\approx 0$ & $\approx 0$ & $11$ quinas interiores, $r = 1$ \\
  \texttt{blocks}    & $0$ & $0$ & $0$ & $0$ & $11$ saltos interiores, $r = 0$ \\
  \texttt{heavisine} & $0$ & $0$ & $16\pi$ & $16\pi$ & $2$ saltos interiores, $r = 0$ \\
  \bottomrule
\end{tabular}
\end{center}

Daí e do Lema~\ref{lem:ordem} sai a tabela que este arquivo existe para
registrar, com $N = 4$ (\texttt{filter.size = 8}, D15):

\begin{center}
\begin{tabular}{lccc}
  \toprule
  & (a) periódico, $\eps = 0$ & (b) margem e extensão & (c) intervalo (CDV) \\
  \midrule
  \texttt{sine}      & $4 = N$ & $4$ & $4$ \\
  \texttt{cosine}    & $4 = N$ & $4$ & $4$ \\
  \texttt{cubic}     & $\mathbf{3/2}$ & $4$ & $4$ (exato) \\
  \texttt{bumps}     & $3/2$ & $3/2$ & $3/2$ \\
  \texttt{blocks}    & $1/2$ & $1/2$ & $1/2$ \\
  \texttt{heavisine} & $1/2$ & $1/2$ & $1/2$ \\
  \midrule
  cenário \texttt{smooth}        & $\mathbf{3/2}$ & $4$ & $4$ \\
  cenário \texttt{inhomogeneous} & $1/2$ & $1/2$ & $1/2$ \\
  \bottomrule
\end{tabular}
\end{center}

O $\seff$ do cenário é o \emph{mínimo} sobre as suas componentes, porque o
viés do Corolário~1 de E1.3 soma sobre os $pq$ blocos e a componente mais
lenta domina a soma.

\subsection*{Por que a cúbica lê $3/2$ num regime e $4$ nos outros dois}

Os coeficientes de \texttt{cubic} foram escolhidos de modo que
$g(0) = g(1) = 0$: $1 - 1{,}4 + 0{,}4 = 0$. A extensão $1$-periódica é
portanto \emph{contínua}, e a Proposição~2 de E1.3, que mede o custo de um
salto de valor, não se aplica. O que não casa é a derivada:
$g'(u) = 3u^{2} - 2{,}8u + 0{,}4$ dá $g'(0) = 0{,}4$ e $g'(1) = 0{,}6$, um
salto de $0{,}2$. Pelo Lema~\ref{lem:ordem} com $r = 1$, a queda por nível é
$r + 1/2 = 3/2$. Esse é o regime~(a), o da teoria, e é a razão do número que
D27 declara.

Nos regimes (b) e (c) a emenda deixa de existir e sobra um polinômio de grau
$3$, que é $C^{\infty}$; pelo teto do Lema~\ref{lem:ordem}, a queda é $N = 4$.
No regime (c) há mais: grau $3 < N = 4$, e a base CDV reproduz exatamente os
polinômios de grau $< N$ em $[0,1]$, de modo que \emph{todos} os coeficientes
de wavelet se anulam e a componente inteira mora em $V_{\jzero}$, na parte não
penalizada. A conferência encontra os coeficientes da cúbica na CDV abaixo do
próprio piso de avaliação da base, em todos os níveis.

\subsection*{Por que \texttt{blocks} e \texttt{heavisine} leem $1/2$ nos três}

Porque o salto é interior, e nenhuma escolha de borda o alcança. Vale
registrar que as três funções de Donoho e Johnstone, como
\texttt{wafc\_component()} as escreve, \emph{também} não saltam na emenda: em
\texttt{blocks} as alturas somam zero
($4-5+3-4+5-4{,}2+2{,}1+4{,}3-3{,}1+2{,}1-4{,}2 = 0$), logo
$g(0) = g(1) = 0$; em \texttt{heavisine} os dois termos de sinal se cancelam
nos dois extremos e o seno tem período $1/4$; e em \texttt{bumps} os núcleos
mais próximos das bordas já decaíram a $10^{-5}$ da amplitude. As seis
componentes do \texttt{dgp.R} são, portanto, todas contínuas na emenda, e a
única que paga alguma coisa por periodização é a cúbica, e paga na derivada.
Isso não é acidente do projeto: são funções desenhadas para testes de wavelets
em $[0,1]$.

Daí a leitura correta do número de D27: o $1/2$ do cenário não homogêneo é
\emph{independente de regime}, e o $3/2$ do cenário suave é
\emph{inteiramente um artefato de periodização}, que sumiria se a teoria
fosse enunciada com margem ou no intervalo.

\subsection*{O que isso muda em E2.4 e em E4}

\begin{enumerate}[label=\textup{(\alph*)}, leftmargin=*]
  \item \textbf{A regra teórica de $J_{n}$ muda de um nível.} Com
    $2^{J_{n}} \ge (n/\log n)^{1/(2\seff+1)}$ e $n = 1000$: $\seff = 3/2$ dá
    $2^{J_{n}}\ge 3{,}47$, isto é $J_{n} = 2$; $\seff = 4$ dá
    $2^{J_{n}}\ge 1{,}74$, isto é $J_{n} = 1$. E2.3 já mediu que a regra fica
    dois níveis \emph{abaixo} do $J$ do oráculo da grade com $\seff = 3/2$;
    declarar $\seff = 4$ a levaria a três. A conclusão de E2.3 (``a regra é
    muito sensível a $\seff$, que não se conhece'') fica mais forte, porque
    agora se vê que $\seff$ depende até da convenção de borda.
  \item \textbf{Há uma discrepância declarada entre o que a teoria diz e o
    que o código gera.} O atributo que D27 manda gravar em \texttt{dgp.R} é o
    número do regime~(a), que é o regime do manuscrito. Mas o piloto roda com
    \texttt{rescale = TRUE} e $\eps = 0{,}05$, isto é, no regime~(b), onde o
    mesmo cenário suave é mais liso. A escolha continua defensável (o atributo
    documenta a teoria enunciada, e a Parte~B mostra que a margem não compra
    taxa), mas tem de ficar escrita, e é matéria do handoff.
  \item \textbf{O teto é $N$, não a função.} Trocar \texttt{filter.size}
    muda o $4$ da tabela: $L = 4$ daria $N = 2$ e o cenário suave leria $2$ em
    (b) e (c), e continuaria em $3/2$ em (a). D15 fixou $L = 8$, e a tabela é
    dele.
\end{enumerate}

\section{O que isto não cobre}

\begin{itemize}[leftmargin=*]
  \item \textbf{A Parte~A é transferência, não reprova.} A
    Proposição~\ref{prop:cdv} aponta, item a item, o passo da prova original
    que sobrevive e por quê; o que se prova aqui de novo é o
    Lema~\ref{lem:pontual} e a aritmética de $\pzero$. Quem levar a rota ao
    manuscrito tem de reescrever as provas de E1.3 a E1.6 com a base trocada,
    o que é trabalho de redação, não de matemática nova.
  \item \textbf{A constante $\Cpsiint$ não é afirmada justa.} O
    Lema~\ref{lem:pontual} dá uma cota, e a conferência mede
    $\approx 7{,}55$ para Daublets de filtro 8; nada aqui diz que a
    dependência em $\jzero$ seja a melhor possível.
  \item \textbf{O limiar de degenerescência da
    Proposição~\ref{prop:cruzamento}(i) é suficiente, não necessário.} A
    conferência encontra a primeira direção nula até um nível antes dele,
    porque uma combinação de funções de escala pode se anular fora da margem
    sem que nenhuma delas caiba nela. A faixa admissível de $J$ é, portanto,
    a medida, e não a do limiar.
  \item \textbf{Nada aqui exibe a constante de compatibilidade sobre as
    direções estimáveis} no caso $\eps > 0$. É a lacuna que E1.3b abriu e que
    D26 decidiu não fechar: a rota do intervalo a contorna em vez de resolvê-la.
  \item \textbf{A ordem exata do Lema~\ref{lem:ordem} é a menos da fase
    diádica.} $\Lambda_{r}(\tau)$ oscila com $\tau = 2^{\lev}t - \tsl$, e uma
    singularidade em ponto diádico pode dar coeficientes anomalamente pequenos
    em alguns níveis. É por isso que as quedas medidas oscilam e que a
    conferência lê médias sobre janelas de vários níveis, e não transições
    isoladas.
  \item \textbf{A Parte~C é sobre as componentes como \texttt{dgp.R} as
    escreve.} Mudar uma constante de \texttt{cubic} para quebrar
    $g(0) = g(1)$ faria o cenário suave cair de $3/2$ para $1/2$ no
    regime~(a).
  \item \textbf{O manuscrito não muda} (D26). Este arquivo é material de
    reserva e de documentação dos cenários.
\end{itemize}

\section{Conferência numérica}

\texttt{derivations/check/07-rota-intervalo.R} (Daublets, \texttt{filter.size
= 8}, $\jzero = 4$ na CDV; grades de $2^{14}$ e $2^{15}$; imprime \texttt{OK}
em cerca de $2$ minutos). O que ela mediu em 2026-09-19:

\paragraph{Parte A.}
\begin{itemize}[leftmargin=*, nosep]
  \item (A1) A base CDV é ortonormal em $L_{2}[0,1]$ com a medida de Lebesgue:
    $\max|G - I|$ vai de $1{,}8\times10^{-6}$ ($J = 4$) a
    $3{,}9\times10^{-4}$ ($J = 8$) na grade de $2^{14}$, que é erro de
    quadratura. É a única propriedade que a parte~(i) da Proposição~3 de E1.4
    usa.
  \item (A2) A reparametrização \eqref{eq:bloco} em $J = 6$: $\mu\T\mu = 1$ a
    $5\times10^{-6}$ (isto é $\norm{\mathbf{1}}_{L_{2}[0,1]} = 1$),
    $\mu\T Q = 0$ a $10^{-10}$, $(\bPhi Q)\T(\bPhi Q) = I$ e $\bPhi Q$
    ortogonal às wavelets a $5\times10^{-3}$, e $15 + 48 = 63 = 2^{6}-1$
    colunas no bloco, como no periódico. Toda coluna tem integral zero.
  \item (A3) $\sup_{u}\norm{W(u)}_{2}^{2}/2^{J}$ vale $7{,}55$, $7{,}55$,
    $7{,}52$, $7{,}42$, $7{,}23$ e $6{,}85$ em $J = 4,\dots,9$ na CDV, contra
    $1{,}30$ a $1{,}36$ na periódica: $\Cpsiint$ é finito e estável, e é
    $5{,}6$ vezes o da base periódica. A reparametrização não aumenta a soma,
    como o Lema~\ref{lem:pontual} prevê ($7{,}546$ contra $7{,}608$ em
    $J = 4$).
  \item (A4) Com $p = 3$, $q = 2$, $\jzero = 4$: $\pzero = 93$ contra $3$;
    $\Amat^{\mathrm{int}}$ tem posto $93$ em $n = 1000$; e
    $\E\normn{\PA\bveps}^{2}$ medido sobre $200$ réplicas é $0{,}0922$ contra
    os $\sigma^{2}\pzero/n = 0{,}0930$ previstos, e $0{,}00295$ contra
    $0{,}00300$ no desenho periódico.
\end{itemize}

\paragraph{Parte B.}
\begin{itemize}[leftmargin=*, nosep]
  \item (B1) Família $\eps = a(L-1)2^{-J}$, $J = 5,\dots,9$. $\lmin(G_{\eps})$:
    $1$ em $a = 0$ (ortonormal), $3{,}5\times10^{-4}$ em $a = 0{,}25$
    (estável em $J$), $1{,}3\times10^{-14}$ em $a = 0{,}5$ e nulo em $a = 1$.
    Queda por nível do erro restrito de $\exp(u)$ em $[\eps, 1-\eps]$:
    $0{,}50$, $0{,}50$, $0{,}08$ e $-0{,}39$. Nenhum $a$ dá ao mesmo tempo
    Gram viva e queda acima de $1/2$ bit, que é a
    Proposição~\ref{prop:cruzamento}(iii) medida; e onde a Gram morre o erro
    restrito desaba de $2{,}3\times10^{-2}$ ($\eps = 0$) a
    $10^{-7}$, que é a Observação~\ref{rem:zero}.
  \item (B2) Com $\eps$ fixo, o maior $J$ com $\lmin(G_{\eps}) > 10^{-8}$ é
    $7$, $5$ e $4$ para $\eps = 0{,}02,\ 0{,}05,\ 0{,}10$, contra o limiar
    suficiente $\lfloor\log_{2}T(\eps)\rfloor = 7$, $6$ e $5$ e contra o
    $\lceil\log_{2}2T(\eps)\rceil = 9$, $8$ e $7$ que excluiria a faixa
    contaminada. As duas faixas não se encontram para nenhum $\eps$.
\end{itemize}

\paragraph{Parte C.}
\begin{itemize}[leftmargin=*, nosep]
  \item (C1) Nenhuma das seis salta na emenda: o salto relativo de $g$ é
    $\le 2{,}5\times10^{-5}$ em todas (o máximo é \texttt{bumps}). O salto
    relativo de $g'$ é $\le 1{,}3\times10^{-5}$ em cinco delas e
    $0{,}334$ na cúbica, que é $0{,}2$ contra $\sup|g'| = 0{,}5996$.
  \item (C2) Queda média por nível de $\norm{\theta_{\lev\cdot}}_{2}$, na
    janela em que a medição é legítima:
    \begin{center}
    \begin{tabular}{lcccccc}
      \toprule
      & \texttt{sine} & \texttt{cosine} & \texttt{cubic} & \texttt{bumps} & \texttt{blocks} & \texttt{heavisine} \\
      \midrule
      (a) periódico $\eps = 0$ & $3{,}99$ & $3{,}95$ & $1{,}45$ & $1{,}37$ & $0{,}48$ & $0{,}51$ \\
      (b) margem e extensão    & $3{,}85$ & $3{,}85$ & $3{,}85$ & $1{,}30$ & $0{,}49$ & $0{,}60$ \\
      (c) intervalo (CDV)      & $5{,}11$ & $4{,}25$ & (no piso) & $1{,}37$ & $0{,}49$ & $0{,}58$ \\
      \bottomrule
    \end{tabular}
    \end{center}
    contra a tabela declarada. A janela começa onde $2^{-\lev}$ fica abaixo da
    menor escala presente (a largura mínima em \texttt{bumps} é $0{,}005$,
    isto é $2^{-7{,}6}$, e no regime (b) entra a largura da transição do corte,
    $\eps/2 = 0{,}05$) e termina no último nível acima de dez vezes o piso de
    avaliação da base, medido pela norma dos coeficientes que \emph{têm} de
    ser exatamente nulos, $\inner{\mathbf{1}}{\psi_{\lev\tsl}}$. Esse piso é
    $10^{-16}$ na base periódica e vai de $7\times10^{-8}$ a
    $3\times10^{-4}$ na CDV, o que é o que impede medir a cúbica ali: os seus
    coeficientes ficam em $10^{-10}$ a $5\times10^{-9}$, abaixo do piso em
    todos os níveis, que é a reprodução exata.
  \item (C3) O mínimo sobre as componentes dá $\seff = 3/2$ e $1/2$ nos dois
    cenários no regime (a), que é D27, e $4$ e $1/2$ nos regimes (b) e (c).
\end{itemize}

\section*{Referências}
\sloppy

Todas as obras citadas aqui já estavam em
\texttt{docs/referencias-verificadas.bib}, conferidas em L1 e L3; E1.8 não
acrescenta nenhuma.

\begin{itemize}[nosep]
  \item Cohen, A. (2003). \emph{Numerical Analysis of Wavelet Methods}.
    Studies in Mathematics and its Applications 32. Elsevier.
    (\texttt{Cohen-2003}.) A \S 3.9, ``Bounded domains'', é a caracterização
    de $\Besov{s}{\pi}{r}[0,1]$ pela base do intervalo que a Parte~A invoca.
  \item Cohen, A., Daubechies, I. e Vial, P. (1993). Wavelets on the interval
    and fast wavelet transforms. \emph{Applied and Computational Harmonic
    Analysis} 1, 54--81. (\texttt{cohen1993wavelets}.) A construção da base,
    as funções de borda e a reprodução de polinômios de grau $< N$.
  \item Daubechies, I. (1992). \emph{Ten Lectures on Wavelets}. CBMS-NSF
    Regional Conference Series in Applied Mathematics 61. SIAM.
    (\texttt{Daubechies-1992}.) A base periodizada é a \S 9.3, p.~304.
  \item Donoho, D. L. e Johnstone, I. M. (1994). Ideal spatial adaptation by
    wavelet shrinkage. \emph{Biometrika} 81, 425--455.
    (\texttt{Donoho-Johnstone-1994}.) As funções \texttt{bumps},
    \texttt{blocks} e \texttt{heavisine} da Parte~C, na forma que
    \texttt{wafc/R/dgp.R} usa.
  \item Montoril, M. H., Pinheiro, A. e Vidakovic, B. (2019). Wavelet-based
    estimators for mixture regression. \emph{Scandinavian Journal of
    Statistics} 46, 215--234.
    (\texttt{Montoril-Pinheiro-Vidakovic-2019}.) O reescalonamento
    $\eps = 1{,}9^{-J}$ herdado pelo \texttt{wall()}, que a
    Observação~\ref{rem:zero} situa logo acima do cruzamento.
\end{itemize}

Arquivos do próprio projeto citados sem repetir a prova:
\texttt{derivations/01-identificabilidade.md} (§5, a reparametrização),
\texttt{02-aproximacao-besov.tex} (Lemas 2 e 3, Corolário 1, Proposição 2,
Lema 10 e as suas observações), \texttt{03-desenho-produtos.tex}
(Proposição 3), \texttt{04-oraculo.tex} (Lemas 4 a 7, Teorema 1) e
\texttt{05-taxas.tex} (Proposição 4, Teorema 2, Corolários 4 e 5).

\end{document}
